\section{Introduction}
\label{sec:introduction}

Variational methods turn differential equations into operator equations on
function spaces. This viewpoint is especially effective for elliptic boundary
value problems: weak solutions require fewer classical derivatives, and the
same bilinear form directly defines a stable numerical method. Standard
references include \cite{ciarlet1978,brenner2008,ern2004}.

Let $\Omega\subset\R^2$ be a bounded polygonal domain. We consider
\begin{equation}
  -\nabla\!\cdot\!\bigl(A(x)\nabla u\bigr)+c(x)u=f
  \quad\text{in }\Omega,
  \qquad
  u=0
  \quad\text{on }\partial\Omega.
  \label{eq:strong-problem}
\end{equation}
The diffusion tensor $A$ is assumed symmetric and uniformly positive definite:
there are constants $0<\alpha\leq\beta$ such that
\begin{equation}
  \alpha\lvert\xi\rvert^2
  \leq \xi^{\mathsf T}A(x)\xi
  \leq \beta\lvert\xi\rvert^2
  \qquad
  \text{for a.e. }x\in\Omega,
  \quad \xi\in\R^2.
  \label{eq:uniform-ellipticity}
\end{equation}
We also assume $c\in L^\infty(\Omega)$ with $c\geq0$ and
$f\in L^2(\Omega)$.

The central question is quantitative: if $u_h$ is computed in a finite
dimensional space with mesh size $h$, how rapidly does $u_h$ approach $u$?
The answer separates into stability and approximation. Galerkin
orthogonality supplies the stability mechanism; interpolation theory supplies
the approximation rate.

\paragraph{Structure of the paper.}
\Cref{sec:weak-formulation} establishes the continuous problem.
\Cref{sec:galerkin} introduces the conforming method, and
\cref{sec:error-estimates} proves energy- and $L^2$-norm estimates.
The experiment in \cref{sec:numerics} checks the asymptotic rates.

% DEMO: Add one sentence here explaining why coercivity is indispensable.
